\documentclass[twocolumn,10pt]{article}
\usepackage[margin=0.75in]{geometry}
\usepackage{amsmath,amssymb,amsthm}
\usepackage{graphicx}
\usepackage{xcolor}
\usepackage[colorlinks=true,linkcolor=blue,citecolor=blue,urlcolor=blue]{hyperref}
\usepackage{authblk}
\usepackage{abstract}
\renewcommand{\Affilfont}{\small}
\setlength{\affilsep}{0.3em}
\newtheorem{theorem}{Theorem}
\newtheorem{lemma}{Lemma}
\newtheorem{proposition}{Proposition}
\newtheorem{definition}{Definition}
\newcommand{\Pp}{\Pi_{\rm probe}}
\newcommand{\R}{\mathcal{R}}
\newcommand{\obs}{f_{\rm obs}}

\title{\vspace{-2em}\textbf{Can cosmology confirm that dark energy is a cosmological constant?
An observable-subspace bound}}
\author[]{Eric Schultz}
\affil[]{Independent researcher \quad ORCID: 0009-0006-6283-1696}
\date{\today}

\begin{document}
\maketitle

\begin{abstract}
\noindent
The simplest hypothesis for dark energy is that it is vacuum energy---a
cosmological constant with equation of state $w=-1$ and no perturbations. We ask
a precise question: given a finite suite of cosmological probes, which
deviations from a cosmological constant are observable even in principle? We
formalise the observability of a dark-energy feature to a probe as the norm of
the feature's linear response projected onto the probe's sensitivity subspace,
$\|\Pi_{\rm probe}\,\mathcal{R}\|$, and show that the three questions answered
separately in a companion programme---the detectability of a localised
background feature, the discrimination of horizon-entropy models, and the
closure of the structure-growth channel---are one construction under three
projections. The finiteness of any real probe suite then implies a sharp
epistemic bound: the observable part of the dark-energy response space is the
span of the probe kernels, and its orthogonal complement is invisible in
principle. Using physically-motivated kernels for six probes, we find that the
cosmological-constant hypothesis \emph{is} falsifiable in the directions the
data cover---low-redshift background evolution and any clustering are well
observed---but that a structurally-defined family of near-constant
alternatives, deviating only at high redshift, on super-horizon scales, or in a
mid-redshift observational desert, is indistinguishable from vacuum energy. The
strongest statement current cosmology can support is therefore ``dark energy is
a cosmological constant within the observable subspace,'' never without
qualification. We map the blind family explicitly.
\end{abstract}

%========================================================================
\section{Introduction}
\label{sec:intro}

The most economical account of cosmic acceleration is that the dark energy
driving it is the energy of the vacuum: a cosmological constant $\Lambda$ with
equation of state $w=-1$ exactly and, being Lorentz invariant, no spatial
perturbations at all. Every current observation is consistent with this
account. It is also the account with the least explanatory content---it names
the substance without addressing why its density takes the value it does---and
so a large literature seeks dynamical alternatives whose deviations from
$w=-1$ might be detected.

This paper asks a question logically prior to that search: setting aside whether
a given experiment has the sensitivity, \emph{which} deviations from a
cosmological constant are observable \emph{even in principle} by a finite suite
of probes? The question is not rhetorical. A cosmological constant is a single
point in an infinite-dimensional space of possible dark-energy behaviours---a
density that could in principle depend on scale and time in arbitrarily many
ways---and a finite set of experiments can constrain only a finite-dimensional
projection of that space. The complement is not merely poorly measured; it is
inaccessible to those probes for any exposure.

We make this precise and find that it has teeth. The observability of a
dark-energy feature reduces to a single object, and the various detectability
and discrimination results established elsewhere in this programme are its
instances. The finiteness of the probe suite then yields an explicit
observable subspace and an explicit blind complement. Applied to the
cosmological-constant hypothesis, the result is neither the complacent ``$\Lambda$
fits, so dark energy is vacuum energy'' nor the sceptical ``we can never know'':
it is a sharp intermediate statement. Vacuum energy is falsifiable in the
directions the data cover, and a specific, mapped family of alternatives is
forever indistinguishable from it.

%========================================================================
\section{The observability functional}
\label{sec:functional}

Let $\mathcal{R}(k,a)$ denote the linear response of a dark-energy model---the
perturbation it induces, as a function of comoving wavenumber $k$ and scale
factor $a$, in the quantities cosmology can measure. A cosmological constant has
$\mathcal{R}\equiv0$ in the perturbed sector and a fixed background; a dynamical
model deviates from this in a pattern $\mathcal{R}(k,a)$ characteristic of its
microphysics.

A probe $P$ does not measure $\mathcal{R}$ pointwise. It measures a weighted
projection: its sensitivity is concentrated on a subspace of the $(k,a)$
response space, described by a kernel $W_P(k,a)$. We work on the response space
$L^2$ of square-integrable functions on the $(\ln k,\ln a)$ domain with inner
product
\begin{equation}
\langle \mathcal{R}_1,\mathcal{R}_2\rangle
=\int \mathcal{R}_1(k,a)\,\mathcal{R}_2(k,a)\,\mu(k,a)\,d\ln k\,d\ln a,
\label{eq:inner}
\end{equation}
and induced norm $\|\mathcal{R}\|^2=\langle\mathcal{R},\mathcal{R}\rangle$. The
weight $\mu(k,a)$ is where survey geometry enters: in a full treatment it is the
Fisher metric, $\mu\propto$ the inverse data covariance folded with the number
of independent modes per cell (the cosmic-volume factor $\sim k^3$ within the
survey footprint). In the schematic evaluation below we take $\mu$ flat on the
logarithmic grid, which is the choice that makes the observable subspace depend
only on the \emph{support} of each probe kernel rather than on assumed noise
levels; we return to this choice as a limitation. We define the observability of
a feature $\mathcal{R}$ to the probe as
\begin{equation}
\obs^{(P)}[\mathcal{R}]
   \;=\; \frac{\|\Pi_P\,\mathcal{R}\|}{\|\mathcal{R}\|},
\label{eq:obsfunc}
\end{equation}
where $\Pi_P$ is the orthogonal projection (in the inner product
Eq.~\eqref{eq:inner}) onto the span of the probe kernel. For a suite of probes
$\{P_i\}$ the relevant projector is onto the combined span,
$\Pi_{\rm probe}=\Pi_{\,{\rm span}\{W_{P_i}\}}$, and $\obs\in[0,1]$ measures the
fraction of the feature's response that lies within reach of the suite.

This is deliberately the information-geometric criterion made concrete: $\obs$
is the cosine between a feature and the observable subspace, and a Fisher
analysis with the same kernels would assign a feature with $\obs\!\to\!0$ an
unbounded confidence interval. The content of what follows is not the
criterion---which is standard---but its \emph{consequences} when the response
space is that of dark energy and the projections are the specific ones the
literature uses.

\paragraph{Three questions, one functional.}
The programme this paper closes answered three questions with what appear to be
three different methods. Each is Eq.~\eqref{eq:obsfunc} under a particular
$\Pi_{\rm probe}$:
\begin{itemize}\itemsep2pt
\item A localised feature in the background expansion (a step or bump in the
dark-energy density) is detectable iff its response is not projected out by the
degeneracy directions of the distance--redshift data---$\Pi$ is the
nuisance-orthogonal Fisher projection onto the space orthogonal to
$\{\Omega_m,\text{distance calibration}\}$.
\item Horizon-entropy models are discriminable at the background level iff their
responses are separated after the same projection---the discrimination question
is Eq.~\eqref{eq:obsfunc} applied to differences of models.
\item The structure-growth channel is closed for event-horizon models because
their perturbation response is annihilated by the sub-horizon projection: the
growth probe's kernel is supported where a null-horizon response scales as
$1/k^3$, so $\Pi_{\rm growth}\mathcal{R}\to0$.
\end{itemize}
The projections genuinely differ; the object projected, and the meaning of the
norm, do not. This is the sense in which the three results are one.

%========================================================================
\section{The observable subspace is finite}
\label{sec:finite}

\begin{proposition}[Observable-subspace bound]
\label{prop:finite}
A suite of $N$ probes with kernels $\{W_{P_i}\}$ renders observable only the
$N$-dimensional (at most) subspace $\mathcal{O}={\rm span}\{W_{P_i}\}$ of the
response space. Any feature $\mathcal{R}\in\mathcal{O}^\perp$ has
$\obs[\mathcal{R}]=0$: it is unobservable to the suite for any exposure.
\end{proposition}

The proposition is immediate---the span of $N$ vectors has dimension at most
$N$---but its force comes from the response space being infinite-dimensional.
No finite programme closes it. Increasing exposure shrinks error bars
\emph{within} $\mathcal{O}$; it does not enlarge $\mathcal{O}$. Only a
genuinely new probe, with a kernel outside the current span, adds a dimension.

The practical version replaces exact orthogonality with smallness. A feature
with $\obs[\mathcal{R}]$ near the value a random response direction would
attain, $\sqrt{N/D}$ for a $D$-dimensional discretisation, is observationally
indistinguishable from unobservable. It is this thresholded, quantitative form
we now apply.

%========================================================================
\section{Can vacuum energy be confirmed?}
\label{sec:lambda}

We discretise the response space on a $(k,a)$ grid spanning
$k\in[10^{-4},1]\,h\,{\rm Mpc}^{-1}$ and $a\in[0.02,1]$ ($D=2400$ cells) and
build physically-motivated kernels for six probes: background distance
(supernovae and baryon acoustic oscillations), the growth rate $f\sigma_8$, the
integrated Sachs--Wolfe effect, weak lensing, primary cosmic microwave
background, and CMB lensing. Each kernel encodes where in $(k,a)$ that probe has
support---low redshift and sub-horizon for growth and lensing, very large scales
for the ISW effect, high redshift for the CMB.

Two cautions must be stated at the outset, because they bound what the numbers
below can mean. First, these kernels are \emph{schematic}: analytic weight
functions capturing the $(k,a)$ support of each probe, not Fisher forecasts from
survey specifications and a Boltzmann code. Second, treating each probe as a
single basis vector \emph{understates} the data: a real tomographic survey
measures many redshift-binned spectra and spans a multi-dimensional subspace,
not one direction. Both push the same way---a realistic suite has a larger
observable subspace $\mathcal{O}$ than our six-vector caricature---so the
specific overlap fractions we report are not to be read as forecasts. We
therefore do not claim the decimal values; we claim the \emph{ordering} they
exhibit, and we test that the ordering is what survives.

It does. Expanding each probe into $n_{\rm bin}$ tomographic slices raises the
subspace dimension (from $6$ to $\sim36$ at $n_{\rm bin}=10$) and lifts every
overlap, but it lifts the low-redshift background direction to near unity while
the blind directions rise far less and stay well below it:
\begin{table}[t]
\caption{Observability $\obs$ of a falsifiable direction (low-$z$ background)
and two blind directions as the observable subspace is enlarged by tomographic
binning. The absolute values rise with dimension $r$, but the ordering
(low-$z$ overlap $\gg$ blind overlap) is preserved throughout. The random
floor $\sqrt{r/D}$ is shown for reference.}
\label{tab:tomography}
\begin{center}\small
\begin{tabular}{cccccc}
\hline\hline
$n_{\rm bin}$ & $r$ & low-$z$ bg & high-$z$ DE & mid-$z$ desert & floor \\ \hline
1  & 6  & 0.85 & 0.21 & 0.22 & 0.05 \\
3  & 14 & 0.93 & 0.31 & 0.35 & 0.08 \\
5  & 20 & 0.95 & 0.49 & 0.34 & 0.09 \\
10 & 36 & 0.96 & 0.52 & 0.37 & 0.12 \\
\hline\hline
\end{tabular}
\end{center}
\end{table}
The gap narrows as the data improve---as it must, since a richer suite sees
more---but it does not close: at every dimensionality the low-redshift
background direction is the most observable and the high-redshift and mid-$z$
directions remain the least, closest to the floor. The qualitative result, a
falsifiable set of directions cleanly separated from a blind set, is a property
of \emph{where the probes have support}, not of the caricature's dimension. What
a full Fisher forecast would change is how far into the blind zones a given
survey generation reaches; what it would not change is that those zones are the
last to be reached.

For definiteness we report the six-vector values below, with the explicit
understanding that they are illustrative markers of the ordering, not forecasts;
a random feature attains $\obs\approx\sqrt{6/2400}\approx0.05$ (the ``blind''
reference), each probe's own kernel returns $\obs=1$, and a smooth physical
response---being concentrated, not white---attains a higher baseline than a
random one, which is why even the blind directions exceed $0.05$ while remaining
far below the observable ones.

A cosmological constant is not unfalsifiable: two ways of departing from it are
well observed (Table~\ref{tab:falsifiable}). The low-redshift background
direction is what current surveys target when they report a preference for
evolving $w$; the clustering direction reflects that a cosmological constant has
\emph{zero} perturbation response, so any dark-energy clustering is a large
signal sitting squarely where growth and lensing probes look. If dark energy
clusters at all, or if its background departs from $w=-1$ at low redshift, we
can see that it is not $\Lambda$.

\begin{table}[t]
\caption{Directions in which the cosmological-constant hypothesis is
falsifiable: deviations from $w=-1$ that a current-generation probe suite
observes well. Values are the six-vector overlaps, illustrative of the ordering
(see text); they rise toward unity under tomographic refinement
(Table~\ref{tab:tomography}).}
\label{tab:falsifiable}
\begin{center}\small
\begin{tabular}{lc}
\hline\hline
Deviation from $\Lambda$ & $\obs$ \\ \hline
Low-$z$ background evolution $w(a)\neq-1$ & $0.92$ \\
Presence of clustering ($c_s^2<1$, any) & $0.87$ \\
\hline\hline
\end{tabular}
\end{center}
\end{table}

The opposite holds for an entire structured family of near-constant deviations
(Table~\ref{tab:blind}).

\begin{table}[t]
\caption{Directions in which the cosmological-constant hypothesis is not
confirmable: near-$\Lambda$ deviations that fall in the probe suite's blind
subspace. Values sit near the random floor and, while they rise under
tomographic refinement, remain the least observable directions at every
dimension (Table~\ref{tab:tomography}).}
\label{tab:blind}
\begin{center}\small
\begin{tabular}{lc}
\hline\hline
Deviation from $\Lambda$ & $\obs$ \\ \hline
High-$z$ / early dark energy ($z\gtrsim4$) & $0.21$ \\
Super-horizon ($k\lesssim10^{-3}$) & $0.45$ \\
Mid-$z$ desert ($k\sim10^{-3}$, $z\sim2$--$4$) & $0.22$ \\
\hline\hline
\end{tabular}
\end{center}
\end{table}

These sit at or near the blind floor. A dynamical dark energy that hugs $w=-1$
at low redshift and departs from vacuum energy only in these zones---early in
time, on the largest scales, or in the redshift gap between the CMB and the
low-redshift probes---produces a response almost entirely in
$\mathcal{O}^\perp$. It is observationally indistinguishable from a cosmological
constant, not because the effect is small but because the probes have no
support where it lives.

\paragraph{The bound.}
The two tables are the result. The cosmological-constant hypothesis is
confirmable only in the observable subspace $\mathcal{O}$, and against the blind
family it is not confirmable at all. The strongest claim the data can support is
\begin{quote}
\emph{dark energy is a cosmological constant within the observable subspace,}
\end{quote}
with the qualification carrying real content: it excludes a specific,
enumerable family of dynamical models that will register in no experiment built
from these six probes. ``Dark energy is vacuum energy,'' stated without that
qualification, is not a conclusion cosmology can reach---not for want of
precision, but for want of directions.

%========================================================================
\section{Discussion}
\label{sec:disc}

The bound is epistemic, not dynamical: it says what is knowable about dark
energy, not what dark energy is. It neither supports nor undermines the vacuum
interpretation. It relocates the question. Confirming that dark energy is
vacuum energy is not a matter of measuring $w$ to more decimal places; it
requires a probe with support in the blind zones---an observable sensitive to
high-redshift or super-horizon dark-energy behaviour---without which a family of
alternatives remains permanently degenerate with $\Lambda$.

Three limitations bound the claim. The probe kernels are schematic and each
probe is treated as few basis vectors rather than a tomographic tower; a
quantitative version requires survey Fisher kernels with redshift binning, which
will raise every overlap and enlarge the observable subspace. As
Table~\ref{tab:tomography} shows, this shifts the fractions substantially but
preserves their ordering---the blind directions remain the least observable---so
we claim the ordering and the existence of the blind family, not the decimal
values. The response model for clustering is a Jeans form, not a Boltzmann
output. And the framework addresses observability, not theoretical prior:
Eq.~\eqref{eq:obsfunc} is the standard information criterion, and the
contribution is its concrete consequence for the $\Lambda$ question, not the
criterion itself.

A word on scope is needed to avoid overstating the bound. The blind subspace is
defined by \emph{where in $(k,a)$ the probes have support}, and that is a
property of the data, not of any dark-energy model class. A deviation from
$w=-1$ localised at high redshift, on super-horizon scales, or in the mid-$z$
desert is hard to observe whatever its microphysical origin---geometric or
scalar-field. So the blind family is not restricted to geometric dark energy;
any dynamical model, including quintessence or k-essence, whose departure from
$\Lambda$ lives in those zones is equally degenerate with a cosmological
constant. What \emph{is} restricted to geometric dark energy is the clean
\emph{response-space bookkeeping} of the earlier sections, where the density
perturbation is inherited from the metric and the projections have the simple
form used in Sec.~\ref{sec:functional}; scalar fields carry an independent
kinetic term and a genuine Lagrangian sound speed, so their response space is
larger and their placement in it requires the full perturbation treatment rather
than the geometric shortcut. The epistemic bound---confirmation only within the
observable subspace---holds for both; only the tidy geometric derivation is
specific to the geometric class.

What survives all three is the shape of the answer. A finite cosmology sees a
finite slice of what dark energy could be. The cosmological constant lives on
the observable side of that slice, which is why it fits; a mapped family of its
imitators lives on the blind side, which is why fitting is not confirming.

\begin{thebibliography}{9}
\bibitem{StepPaper} E. Schultz, ``Detectability and exclusion of localized
dark-energy features in the expansion history,'' companion paper (2026).
\bibitem{CompanionPaper} E. Schultz, ``Which generalized entropy?
Discrimination among horizon-entropy dark-energy models and the closure of the
structure-growth channel,'' companion paper (2026).
\bibitem{NullHorizon} E. Schultz, ``Which dark energy can structure growth see?
A response-kernel classification of geometric dark energy,'' companion paper
(2026).
\end{thebibliography}

\end{document}
